\documentclass[11pt,a4paper]{article}

\usepackage[margin=2.4cm]{geometry}
\usepackage[T1]{fontenc}
\usepackage[utf8]{inputenc}
\usepackage{lmodern}
\usepackage{microtype}
\usepackage{amsmath}
\usepackage{amssymb}
\usepackage{booktabs}
\usepackage{longtable}
\usepackage{array}
\usepackage{tabularx}
\usepackage{listings}
\usepackage{xcolor}
\usepackage{hyperref}
\usepackage{enumitem}
\usepackage{titlesec}

\hypersetup{
  colorlinks=true,
  linkcolor=blue!50!black,
  urlcolor=blue!60!black,
  citecolor=black,
  pdftitle={SSLA-RT Engineering Handbook},
  pdfauthor={ssla-rt},
}

\definecolor{codebg}{rgb}{0.97,0.97,0.97}
\definecolor{codecomment}{rgb}{0.4,0.4,0.4}
\definecolor{codekw}{rgb}{0.1,0.3,0.7}
\definecolor{codestr}{rgb}{0.6,0.3,0.1}

\lstdefinestyle{code}{
  backgroundcolor=\color{codebg},
  commentstyle=\color{codecomment}\itshape,
  keywordstyle=\color{codekw}\bfseries,
  stringstyle=\color{codestr},
  basicstyle=\ttfamily\footnotesize,
  breakatwhitespace=false,
  breaklines=true,
  captionpos=b,
  keepspaces=true,
  showspaces=false,
  showstringspaces=false,
  showtabs=false,
  tabsize=2,
  frame=single,
  framerule=0pt,
  framexleftmargin=4pt,
  framexrightmargin=4pt,
  framextopmargin=4pt,
  framexbottommargin=4pt,
  xleftmargin=4pt,
  xrightmargin=4pt,
}
\lstset{style=code}
\lstdefinelanguage{CUDA}{
  morekeywords={__global__,__device__,__shared__,__syncthreads,__syncwarp,
    __launch_bounds__,extern,const,float,int,unsigned,long,bool,void,
    constexpr,if,else,for,while,return,struct,template,class,inline,
    static,volatile,nullptr,true,false,sizeof},
  morecomment=[l]{//},
  morecomment=[s]{/*}{*/},
  morestring=[b]"
}
\lstdefinelanguage{Bash}{
  morekeywords={cmake,make,python3,python,cd,mkdir,sudo,apt,bash,echo,export,for,do,done,if,then,fi,git},
  morecomment=[l]{\#},
  morestring=[b]"
}

\titleformat{\section}{\normalfont\Large\bfseries\color{blue!40!black}}{\thesection}{1em}{}
\titleformat{\subsection}{\normalfont\large\bfseries\color{blue!50!black}}{\thesubsection}{1em}{}
\titleformat{\subsubsection}{\normalfont\normalsize\bfseries}{\thesubsubsection}{1em}{}

\setlength{\parskip}{0.45em}
\setlength{\parindent}{0pt}

% ----------------------------------------------------------------------
\begin{document}

\begin{titlepage}
\vspace*{4cm}
\begin{center}
\Huge\bfseries SSLA-RT Engineering Handbook
\par\vspace{1.5cm}
\Large
Real-Time Event-Camera Inference\\on Jetson Orin NX
\par\vspace{2.5cm}
\large
A consolidated technical record of the design decisions, optimizations,\\
failed experiments, and integration work performed during the\\
porting and tuning of SSLA-S to the Orin NX + DVXplorer Micro target.
\par\vspace{4cm}
\normalsize
2026-05-14\\[0.4cm]
\texttt{ssla-rt} repository --- \texttt{docs/handbook/}
\end{center}
\end{titlepage}

\tableofcontents
\clearpage

% ====================================================================
\section{Scope and Purpose}

This handbook is a single-document reference covering the engineering work
that took the SSLA-S event-camera detector from a research codebase to a
real-time inference runtime on the NVIDIA Jetson Orin NX edge platform.
It is intended to be readable cold by an engineer who must:

\begin{itemize}[itemsep=0pt]
  \item resume the optimization work after a break,
  \item port the runtime to a different Jetson SKU or future GPU,
  \item swap in a different event camera,
  \item adapt the design to a different SSLA backbone variant.
\end{itemize}

It complements three companion documents that already live in the
repository:

\begin{itemize}[itemsep=0pt]
  \item \texttt{docs/PIPELINE\_DESIGN.md} --- the canonical pipeline diagram
        and constraint set.
  \item \texttt{docs/PROFILE.md} --- the authoritative latency and
        throughput numbers (refreshed measurement-by-measurement).
  \item \texttt{docs/DEV\_GUIDE.md} --- the methodological playbook,
        hardware-adaptation checklists, and the running record of
        dead-end optimization attempts.
\end{itemize}

The handbook does not duplicate those files. Where they go into
file-and-line detail, this document presents the consolidated
\emph{narrative} of how the system was built and why it looks the way it
does, so that future readers can recover not only the answer but also
the chain of reasoning that produced it.

% ====================================================================
\section{System at a Glance}

\subsection{Reference platform}

\begin{tabularx}{\linewidth}{lX}
\toprule
Compute & NVIDIA Jetson Orin NX (8\,SMs, sm\_87, SM clock pinned at
          918\,MHz in MAXN). \\
Sensor  & iniVation DVXplorer Micro, 640$\times$480 sensor, USB\,2.0,
          accessed via \texttt{dv-processing} 2.0.3+. \\
Network & SSLA-S: 4-stage recurrent event backbone with a YOLOX-style
          detection head. 0.5\,M parameters, $\sim$ 96 hidden channels
          at the final stage. \\
Throughput target & 1\,Mev/s (events per second), end-to-end. \\
Latency target    & sub-millisecond emit-to-prediction for output-producing
                    events. \\
\bottomrule
\end{tabularx}

\subsection{Pipeline shape}

The deployment splits SSLA-S across CPU and GPU. The CPU handles the
first two backbone stages (layers \texttt{L0}--\texttt{L3}) and the
intermediate spatial downsamples and temporal dropouts. Events are then
pushed to a per-block pinned-host ring; a persistent GPU kernel
consumes that ring, runs stages~3 and~4 (layers \texttt{L4}--\texttt{L7}),
and writes per-cell predictions to a second pinned-host buffer.

\begin{center}
\small
\begin{tabular}{l}
\toprule
\textbf{End-to-end flow} \\
\midrule
\texttt{emit (camera/synth)} \\
\;\;\;\(\downarrow\) \\
\texttt{CPU shard:\;L0\(\to\)L1\(\to\)pool\(\to\)tdrop\(_{s2}\)\(\to\)L2\(\to\)L3\(\to\)pool\(\to\)tdrop\(_{s3}\)} \\
\;\;\;\(\downarrow\)\;\; \emph{push} \\
\texttt{Pinned ring (per-block SPSC, 64 KiB records)} \\
\;\;\;\(\downarrow\)\;\; \emph{pop} \\
\texttt{GPU persistent kernel:\;tdrop\(_{s2}\)\(\to\)L4\(\to\)L5\(\to\)pool\(\to\)tdrop\(_{s3}\)\(\to\)L6\(\to\)L7\(\to\)head} \\
\;\;\;\(\downarrow\) \\
\texttt{Pinned predictions buffer (per-cell, version-tagged)} \\
\bottomrule
\end{tabular}
\end{center}

\subsection{Production performance ceiling}

\begin{tabularx}{\linewidth}{lXr}
\toprule
\textbf{Metric} & \textbf{Setting} & \textbf{Value} \\
\midrule
CPU admit ceiling (stable)   & 8 shards (n=6 sustainable), halo=2, fused NEON layer & 1.85--2.0 Mev/s \\
GPU drain ceiling            & 8 blocks, persistent kernel, balanced 2W$\times$4H topology & 227--230 kev/s \\
Per-event GPU kernel p50     & at saturation                                              & $\sim$310 \(\mu\)s \\
End-to-end emit\(\to\)done p50 & at saturation (output-producing events only)             & $\sim$330 \(\mu\)s \\
\bottomrule
\end{tabularx}

Numbers are the current values at the end of the May 2026 optimization
campaign, with stub weights. They are reproduced from
\texttt{docs/PROFILE.md}\,\S\S 1--2; see there for sweeps, percentiles,
and historical comparison.

% ====================================================================
\section{Hardware-Imposed Invariants}

These constraints are not negotiable design choices: they come from
the physical characteristics of the Orin NX and the SSLA-S architecture.
Violating any of them silently breaks correctness or causes deadlocks.

\subsection{Tegra coherence: pinned memory for cross-CPU/GPU traffic}

Orin NX reports \texttt{CONCURRENT\_MANAGED\_ACCESS = 0}. Practical
consequence: CUDA managed memory is \emph{not safe} to use for data that
the CPU writes while the GPU is concurrently reading. The driver runs
in the legacy migrate-on-launch mode, in which a managed page must be
fully owned by one side at a time.

Every cross-domain buffer therefore lives in \textbf{pinned host memory},
allocated via \texttt{cuMemHostAlloc}:

\begin{itemize}[itemsep=0pt]
  \item per-block ring buffers (CPU writes, GPU polls),
  \item ring head and tail atomics,
  \item the stop flag,
  \item the per-cell predictions buffer.
\end{itemize}

Managed memory (\texttt{cuMemAllocManaged}) is reserved for GPU-only state:

\begin{itemize}[itemsep=0pt]
  \item hidden state (read and written by GPU only),
  \item layer and head weights (CPU writes once at init, GPU reads forever),
  \item tdrop counters (GPU only after init zero),
  \item GPU timing slots.
\end{itemize}

The helpers \texttt{cuda\_util.alloc\_pinned} and
\texttt{cuda\_util.alloc\_managed} (Python) document which class each
buffer should be in. Using the wrong class produces stale reads that
look like correctness bugs.

\subsection{Halo $=2$ at every CPU sharding boundary --- LOCKED}

The CPU pipeline shards in the spatial domain. Two adjacent shards must
overlap by a halo of at least 2 cells at the resolution of the
\emph{coarser} of the two pooling steps that bracket the layer in
question. For SSLA-S, this works out to \textbf{halo = 2 at the s2
grid}, which after the next s1\(\to\)s2 pool still gives the needed
$\pm 1$ at s3 for the 3$\times$3 patch.

Halo = 1 is \emph{not} sufficient because:

\begin{itemize}[itemsep=0pt]
  \item the 3$\times$3 spatial convolution at the current resolution
        needs $\pm 1$ cell of halo at that resolution, but
  \item after the next pool, the same convolution at the next
        resolution needs $\pm 1$ at that resolution, which is $\pm 2$
        cells at the current resolution.
\end{itemize}

This rule is enforced in \texttt{lib\_stage01\_to\_gpu.cpp} and must be
preserved by every future sharding scheme. It is one of the few
explicit \emph{locked} constraints of the design --- the project's
durable rule.

\subsection{Cell-owner warp partition: race-free without GPU atomics}

The production GPU kernel uses 9 warps per CUDA block, indexed
\texttt{warp\_id $\in [0,9)$}. Each warp owns the hidden-state cells
whose 2-D index satisfies
\[
  (\mathrm{cy} \bmod 3,\;\mathrm{cx} \bmod 3) = (\lfloor \mathrm{warp\_id}/3 \rfloor,\;\mathrm{warp\_id} \bmod 3).
\]
For any event whose 3$\times$3 patch needs to update 9 hidden cells,
the 9 cells map to 9 distinct \texttt{(cy\,\%\,3, cx\,\%\,3)} pairs,
hence to 9 distinct owning warps. Therefore \textbf{any given hidden
cell is written by exactly one warp ever}, eliminating the need for
\texttt{atomicAdd} or any cross-warp synchronization on hidden state.

The kernel uses \texttt{\_\_syncthreads()} only between coarse-grained
phases (\textsc{DISPATCH}, \textsc{COMPUTE}, \textsc{GATHER},
\textsc{TDROP}), not inside the per-event step. Within a warp's queue
of (event, delta) tasks, the sequential processing uses only
\texttt{\_\_syncwarp()} and warp shuffles.

If a future kernel restructuring reduces the warp count or changes the
ownership function, the race-free property must be re-proven from
scratch, and the P1 oracle drift gate
(\texttt{bench\_s2\_s3\_head\_celled.py}) will catch any violation
within a few hundred events.

% ====================================================================
\section{CPU Pipeline}

\subsection{Topology}

The CPU side consists of $N$ shard threads (currently $N=6$ on cores
1--6) each handling a contiguous Y-band of the input grid with a 2-cell
overlap to the next shard. Events whose $y$ coordinate falls inside the
band are processed; events on the overlap are processed by both
neighbours and reconciled at the GPU side via ownership (each event
ends up routed to exactly one GPU block by the same cell-owner
function).

Per-shard work for one event:

\[
  \texttt{stage\_0}\;(L_0 + L_1) \;\to\; \texttt{pool\,/2} \;\to\; \texttt{tdrop}_{s2}
  \;\to\; \texttt{stage\_1}\;(L_2 + L_3) \;\to\; \texttt{pool\,/2} \;\to\; \texttt{tdrop}_{s3}
\]

Each stage applies a \emph{layer\_pair}: \texttt{matvec\(\to\)LRU
step\(\to\)matvec\_accum\(\to\)gather\(\to\)LayerNorm}. With halo,
each shard touches roughly $1/N + 2\Delta$ of the grid where $\Delta$
is the halo width in cells.

\subsection{NEON optimization rounds}

The CPU per-event budget went through three optimization rounds, each
documented in \texttt{docs/PROFILE.md} and \texttt{docs/DEV\_LOG.md}.

\subsubsection*{Round 1 (baseline)}

GCC autovectorization on the matvec and LRU loops. Per-shard rate
$\sim$135 kev/s; system ceiling $\sim$0.82 Mev/s at $n=8$ shards.

\subsubsection*{Round 2 (matvec + LRU NEON)}

Hand-written NEON specializations of \texttt{matvec\_ct<12,36>} and
\texttt{matvec\_ct<24,72>} (the IN=12 and IN=24 layers) using 4-output
ILP, plus a NEON sigmoid approximation for the LRU gate. Per-shard rate
rose to 216 kev/s; system ceiling 1.30 Mev/s at $n=7$.

\subsubsection*{Round 3 (fused interior layer kernel)}

A hand-fused inner kernel that holds \texttt{qvg} and \texttt{qh} in
named NEON registers (\texttt{float32x4\_t} variables) across the
matvec\,\(\to\)\,LRU\,\(\to\)\,matvec\_accum sequence, instead of
spilling them to a \texttt{qvg[18]} array between primitives. Drops the
spill cost and lets the compiler keep them register-resident across
the LRU step. Per-shard rate 327 kev/s at $n=4$; system ceiling
\textbf{1.85 Mev/s stable at $n=6$}, with burst $\sim$2.0 Mev/s at
$n=8$.

\subsubsection*{Pitfall worth remembering}

The fused kernel only realized its gain when written with \emph{named}
\texttt{float32x4\_t} variables. A version using
\texttt{float32x4\_t qvg[18]} (array of 18 NEON registers) compiles
correctly but does not deliver the speedup --- the array form pessimizes
the register allocator into spilling. The named-variable form is the
canonical recipe, recorded in \texttt{ssla\_kernels.cpp}.

\subsection{Per-event stage breakdown}

\begin{tabularx}{\linewidth}{lXrrr}
\toprule
& & Round\,1 & Round\,2 & Round\,3 \\
\midrule
\(L_0+L_1\) (\texttt{stage\_forward(0)}) & per event & 3.94 \(\mu\)s & 1.83 \(\mu\)s & \textbf{1.05 \(\mu\)s} \\
\(L_2+L_3\) (\texttt{stage\_forward(1)}) & per event & 12.88 \(\mu\)s & 9.50 \(\mu\)s & \textbf{5.34 \(\mu\)s} \\
Per-shard rate                          & sustained  & 135 kev/s   & 216 kev/s  & \textbf{327 kev/s} \\
System ceiling (stable)                  & best $n$  & 0.82 Mev/s ($n=8$) & 1.30 Mev/s ($n=7$) & \textbf{1.85 Mev/s ($n=6$)} \\
\bottomrule
\end{tabularx}

% ====================================================================
\section{GPU Kernel: \texttt{ssla\_s2\_s3\_head\_celled.cuh}}

This is the production GPU kernel. It is launched once at startup as a
\emph{persistent} kernel and runs until the host sets a stop flag.
There is one CUDA block per GPU block in the topology (currently 8),
each pinned to one SM by virtue of the grid dimension equaling the SM
count.

\subsection{Per-block run loop}

\begin{enumerate}[itemsep=0pt]
  \item Sleep-poll the ring tail until at least one event is published.
  \item Take the next contiguous run of up to \texttt{BATCH = 8} events
        (the batch never spans non-contiguous slots; consumed in order).
  \item Snapshot the batch into shared memory (\texttt{event\_slots}).
  \item Run \texttt{tdrop\(_{s2}\)} for the batch, writing per-event
        \texttt{pass2} flags.
  \item Run \(L_4\) (\texttt{run\_layer\_celled\_implicit}) and \(L_5\)
        (\texttt{run\_layer\_celled\_implicit\_split} with the
        \texttt{pass2} output mask).
  \item Pool s2 coordinates to s3 (right shift by 1).
  \item Run \texttt{tdrop\(_{s3}\)} with the \texttt{pass2} mask,
        writing \texttt{pass3} flags.
  \item Run \(L_6\) and \(L_7\) (same pattern, with \texttt{pass2} as
        enter-mask and \texttt{pass3} as output-mask).
  \item Run the YOLOX head for events with
        \texttt{is\_owner $\wedge$ pass2 $\wedge$ pass3}; write
        per-cell predictions and increment the per-cell version
        counter.
  \item Record per-event timing in the timing-slot ring.
  \item Advance ring tail; update \texttt{events\_done} counter.
\end{enumerate}

\subsection{Implicit-dispatch layer step}

For each layer, the kernel computes, in lock-step across the 9 warps:

\begin{enumerate}[itemsep=0pt]
  \item \textbf{DISPATCH}: each warp scans the batch and finds events
        whose 3$\times$3 patch covers a cell it owns. The corresponding
        \texttt{(event, delta)} pairs go into a per-warp task queue.
  \item \textbf{COMPUTE}: each warp processes its queued tasks
        sequentially: \texttt{qvg matvec\(\to\)LRU step\(\to\)goW
        matvec} into a per-event contribution slot. Hidden state
        writes are exclusive to the owning warp.
  \item \textbf{GATHER}: one warp per event sums the 9 warp
        contributions plus the residual, applies LayerNorm, and writes
        the result back into the event's broadcast slot for the next
        layer.
\end{enumerate}

Cross-warp synchronization is one \texttt{\_\_syncthreads()} between
the three phases. Inside a phase, only \texttt{\_\_syncwarp()} is
needed.

\subsection{Optimizations applied to the kernel}

\subsubsection*{L4 \texttt{qvgIn} shared-memory cache (121 KB)}

The \(L_4\) \texttt{qvgIn} tensor (9 patches $\times$ 3 projections $\times$
$48 \times 24$) is exactly 121 KB. Loading it into shared memory once
per block at kernel entry, and reading from smem rather than from
global memory in the layer-step matvecs, eliminated the L2/DRAM
re-fetch on every batch and removed a measurable L2 thrash signature in
\texttt{nvprof}. Required \texttt{\_\_launch\_bounds\_\_(288, 1)} and
\texttt{cuFuncSetAttribute(MAX\_DYNAMIC\_SHARED\_SIZE\_BYTES)} since
Orin NX sm\_87 supports up to $\sim$167 KB dynamic smem per block.

Realized gain: \(+\) 4\% throughput. (Initial projection was 20--30\% ---
see \S \ref{sec:methodology} for the cost of trusting projections.)

\subsubsection*{L7 \texttt{qvgIn} L2-persistence window (perf bench only)}

\(L_7\)'s \texttt{qvgIn} (972 KB) is too large for smem but fits the L2
working set if pinned via
\texttt{CU\_STREAM\_ATTRIBUTE\_ACCESS\_POLICY\_WINDOW}. This gives
$+1.7$\% offline and approximately $0$\% live. Currently enabled only
in the offline perf bench --- not worth a kernel-level dependency.

\subsubsection*{Balanced 2W $\times$ 4H topology (corner blocks taller)}

The 8-block topology splits the s2 grid into 2 columns and 4 rows. The
default equal-rows split leaves the middle blocks bottlenecked because
mid-image typically contains the most motion. Balancing the Y splits to
5/3/3/5 cells brought middle blocks down to the same rate as corner
blocks ($\sim$28 kev/s each), realizing $+$ 4\% over the equal split.

\subsubsection*{State-only (VG-only) path for tdrop-dropped events at L5/L7}

Events that fail \texttt{tdrop\(_{s2}\)} or \texttt{tdrop\(_{s3}\)}
still need their hidden state to advance (the LRU recurrence is
unconditional), but their final output (Q matrix and post-LN
projection) is unused. The VG-only branch in
\texttt{process\_patch\_cell\_state\_only} skips the Q matvec and the
post-LRU \texttt{qh} compute. Realized gain: $+$ 3\% throughput.

\subsubsection*{A1 optimization: skip head for dropped events}

Events that fail \texttt{tdrop\(_{s3}\)} have their predictions
discarded anyway. The head section skips them entirely (no head
matvec, no preds write, no version bump). Trivial save but
non-negligible at saturation.

\subsection{Smem budget at the production layout}

\begin{tabularx}{\linewidth}{lr}
\toprule
\textbf{Buffer} & \textbf{Bytes} \\
\midrule
event\_slots[BATCH=8] $\times$ (in\_feat + meta)             & $\sim$5\,400 \\
contrib[BATCH][9 warps][OUT\_MAX=96]                          & 27\,648 \\
per\_warp\_in\_feat\_sm[9][96]                                & 3\,456 \\
per\_warp\_qh\_sm[9][96] \emph{(reused for head)}             & 3\,456 \\
task\_event / task\_delta / task\_count / pass2 / pass3      & $\sim$700 \\
L4 qvgIn cache (9 $\times$ 3 $\times$ 48 $\times$ 24)         & 124\,416 \\
\midrule
Total                                                         & $\sim$166\,128 (rounded to 16 B) \\
\bottomrule
\end{tabularx}

% ====================================================================
\section{Real-Time Demo Pipeline}

The live runner (\texttt{orin/hybrid\_runner\_multi.py}) wraps the
production kernel with a DVXplorer camera reader thread, a display
thread (cv2 window), and a stats loop. This section describes the
demo's structure and the pitfalls encountered building it.

\subsection{Camera reader thread}

\begin{lstlisting}[language=Python,caption={Camera reader, simplified.}]
class CameraReader(threading.Thread):
    SS_LADDER = {"EVERY_PIXEL": 1, "EVERY_SECOND": 2,
                 "EVERY_FOURTH": 4, "EVERY_EIGHTH": 8}

    def run(self):
        import dv_processing as dv
        cam = dv.io.camera.open()
        cam.setSubSampleHorizontal(SS_VAL); cam.setSubSampleVertical(SS_VAL)
        cam.setContrastThresholdOn(contrast)
        cam.setContrastThresholdOff(contrast)
        cam.setCropArea((roi_x, roi_y, roi_w, roi_h))

        while not self.stop_event.is_set():
            batch = cam.getNextEventBatch()
            ev = batch.numpy()
            packed = np.empty((ev.size, 4), dtype=np.float32)
            packed[:, 0] = (ev["timestamp"] - t0).astype(np.float32)
            # Translate to ROI origin then subsample.
            packed[:, 1] = ((ev["x"] - roi_x) // ds).astype(np.float32)
            packed[:, 2] = ((ev["y"] - roi_y) // ds).astype(np.float32)
            packed[:, 3] = ev["polarity"].astype(np.float32)
            self.api.submit(packed[valid])
\end{lstlisting}

Three subtle details:

\begin{description}[itemsep=0pt,leftmargin=*]
  \item[Library choice.] iniVation's older C library
        \texttt{libcaer} (3.3.17) cannot open the DVXplorer
        \emph{Micro} variant; only the C++ \texttt{dv-processing}
        library $\geq$ 2.0.3 does. The Python bindings of
        \texttt{dv-processing} are the easiest path on Orin.
  \item[Contrast threshold range.] The DVXplorer Micro chip silently
        clamps the contrast threshold at 17. The
        \texttt{getContrastThreshold()} getter returns whatever value
        was set, regardless of the clamp, so readback does \emph{not}
        reveal that values above 17 are ineffective. The full valid
        range is therefore 0--17, not 0--255 as the API docs suggest.
        Use 6--10 for balanced sensitivity, 2--4 for high sensitivity,
        12--17 for low.
  \item[ROI translation.] \texttt{setCropArea} crops the sensor
        in-hardware, but the event coordinates still arrive in
        sensor-space. Subtract the crop origin before feeding to the
        SSLA grid, otherwise events miss the grid.
\end{description}

\subsection{Display thread}

\subsubsection*{Render at native resolution, scale once}

Naively rendering events at the upscaled display resolution and
splatting scale$\times$scale squares per event costs $O(scale^2)$ per
event and dominates frame time at 30\,fps. The correct pattern is:

\begin{enumerate}[itemsep=0pt]
  \item Build a small canvas of size $H\times W$ (native event resolution).
  \item Splat events at native resolution (one pixel per event).
  \item \texttt{cv2.resize(small, (W$\times$s, H$\times$s),
        INTER\_NEAREST)} for the upscale.
\end{enumerate}

\texttt{cv2.resize} is heavily optimized C, while a Python-level splat
loop pays the GIL repeatedly. Measured difference: from "feels laggy"
to smooth 30\,fps.

\subsubsection*{Frame pacing via \texttt{cv2.waitKey}}

Use \texttt{cv2.waitKey(wait\_ms)} for both the GUI event pump and
frame pacing, not a Python \texttt{time.sleep} loop. The \texttt{waitKey}
call pumps the X11/Wayland event loop, which a sleep does not, and is
also the only portable way to detect the ESC key. Compute the wait
budget as $\max(1, \text{period}_{\text{ms}} - \text{elapsed}_{\text{ms}})$.

\subsubsection*{Non-maximum suppression with \texttt{cv2.dnn.NMSBoxes}}

The per-cell head produces one bounding box per s3 cell. To collapse
duplicates we run NMS across all blocks' candidates:

\begin{lstlisting}[language=Python]
keep_idx = cv2.dnn.NMSBoxes(boxes_xywh, scores,
                            score_threshold=obj_thresh,
                            nms_threshold=nms_iou)
\end{lstlisting}

\textbf{Footgun:} \texttt{cv2.dnn.NMSBoxes} asserts
\texttt{score\_threshold} $\geq$ 0. If the threshold is computed
dynamically (e.g.\ to accommodate a top-K fallback path that includes
sub-threshold candidates), clamp it to zero:
\texttt{eff\_conf = max(0.0, $\dots$)}.

% ====================================================================
\section{Trained-Checkpoint Loading}

The May 2026 work added the ability to run the deploy pipeline with a
real DSEC-trained SSLA-S checkpoint, not just random/stub weights.
This required both Python-side wiring and one structural change to the
GPU kernel (a real 2-layer YOLOX head). The section is detailed
because the integration touched several seemingly unrelated parts of
the code at once.

\subsection{Checkpoint inspection}

The upstream checkpoint is a PyTorch Lightning \texttt{.ckpt} produced
by the openeva research repository. The relevant parts:

\begin{lstlisting}[language=Python,caption={Inspecting the checkpoint.}]
ck = torch.load("best_epoch=41-step=52500.ckpt", map_location="cpu")
sd  = ck["state_dict"]           # live optimizer state
ema = ck["ema_state_dict"]       # EMA-of-weights (decay 0.9999), eval-preferred
hp  = ck["hyper_parameters"]     # dataset (DSEC), height/width, stages, ...
\end{lstlisting}

\textbf{Use the EMA state dict at inference}; the live optimizer state
is generally slightly worse-generalizing. Both have identical
structure for this model.

The state dict contains 96 tensors in two groups: backbone layers (the
8 SSLA-S layers) and a YOLOX-style detection head.

\subsection{Layer index remapping}

The upstream model uses non-contiguous layer indices because the
\texttt{spatial\_downsample} and \texttt{temporal\_dropout} operators
also take index slots in its \texttt{nn.Sequential}-like layout. The
deploy uses consecutive indices 0--7. The conversion table is:

\begin{tabularx}{\linewidth}{lXl}
\toprule
\textbf{Upstream index} & \textbf{Description}                       & \textbf{Deploy index} \\
\midrule
0  & stage 0, layer 0 (K=1, 2 $\to$ 12)                              & 0 \\
1  & stage 0, layer 1 (K=3, 12 $\to$ 12)                              & 1 \\
4  & stage 1, layer 0 (K=3, 12 $\to$ 24)                              & 2 \\
5  & stage 1, layer 1 (K=3, 24 $\to$ 24)                              & 3 \\
8  & stage 2, layer 0 (K=3, 24 $\to$ 48)                              & 4 \\
9  & stage 2, layer 1 (K=3, 48 $\to$ 48)                              & 5 \\
12 & stage 3, layer 0 (K=3, 48 $\to$ 96)                              & 6 \\
13 & stage 3, layer 1 (K=3, 96 $\to$ 96)                              & 7 \\
\bottomrule
\end{tabularx}

Indices 2,3,6,7,10,11 are stateless ops (pool, tdrop) with no
parameters and so are not exported.

The conversion script \texttt{tools/convert\_dsec\_ckpt.py} performs
this remap, drops unused buffers
(\texttt{pixel2patch\_lut}, \texttt{num\_batches\_tracked}), and writes
the standard \texttt{weights.npz + meta.json} pair that
\texttt{libstage01\_to\_gpu} expects.

\subsection{2-layer YOLOX head structure}

The trained head is markedly more complex than the random
single-linear head used during early stub-only testing:

\begin{tabularx}{\linewidth}{lXl}
\toprule
\textbf{Tensor} & \textbf{Shape} & \textbf{Role} \\
\midrule
\texttt{head/cls\_convs/0/0/weight} & (96, 96, 1, 1)             & 1$\times$1 conv, classification branch \\
\texttt{head/cls\_convs/0/0/bias}   & (96,)                       & conv bias (absent in stub schema) \\
\texttt{head/cls\_convs/0/1/*}      & (96,)$\times$4              & BatchNorm (\(\gamma, \beta, \mu, \sigma^2\)) \\
\texttt{head/cls\_preds/0/weight}   & (2, 96, 1, 1)              & 2 class logits per cell \\
\texttt{head/cls\_preds/0/bias}     & (2,)                       & \\
\texttt{head/reg\_convs/0/*}        & (96, 96, ...)              & 1$\times$1 conv, regression/obj branch \\
\texttt{head/reg\_preds/0/weight}   & (4, 96, 1, 1)              & 4 bbox parameters (tx, ty, log\_w, log\_h) \\
\texttt{head/reg\_preds/0/bias}     & (4,)                       & \\
\texttt{head/obj\_preds/0/weight}   & (1, 96, 1, 1)              & objectness logit \\
\texttt{head/obj\_preds/0/bias}     & (1,)                       & \\
\bottomrule
\end{tabularx}

Per-cell forward pass:

\begin{align*}
y_{\text{cls}} &= \mathrm{SiLU}(W_{\text{cls\_conv}} \cdot x + b_{\text{cls\_conv}}) \\
\text{cls\_logits} &= W_{\text{cls\_pred}} \cdot y_{\text{cls}} + b_{\text{cls\_pred}} \\[0.3em]
y_{\text{reg}} &= \mathrm{SiLU}(W_{\text{reg\_conv}} \cdot x + b_{\text{reg\_conv}}) \\
\text{reg\_out} &= W_{\text{reg\_pred}} \cdot y_{\text{reg}} + b_{\text{reg\_pred}} \\
\text{obj\_logit} &= W_{\text{obj\_pred}} \cdot y_{\text{reg}} + b_{\text{obj\_pred}}
\end{align*}

The two branches share the input $x$ (the L7 output, $C_3 = 96$
channels) but have \emph{separate} 96$\to$96 conv layers. The
non-linearity SiLU prevents fusing into a single matrix.

\subsection{Folding BatchNorm into the conv weights}

BatchNorm at eval mode is a per-channel affine transform:
\[
y_{\text{bn}}[c] = \gamma[c] \cdot \frac{y[c] - \mu[c]}{\sqrt{\sigma^2[c] + \epsilon}} + \beta[c]
                  = \mathrm{scale}[c] \cdot y[c] + \mathrm{shift}[c]
\]
with $\mathrm{scale} = \gamma / \sqrt{\sigma^2 + \epsilon}$ and
$\mathrm{shift} = \beta - \mathrm{scale} \cdot \mu$. Composing with the
preceding 1$\times$1 conv:
\[
W^{\text{fused}}_{c_{\text{out}}, c_{\text{in}}} = \mathrm{scale}[c_{\text{out}}] \cdot W^{\text{conv}}_{c_{\text{out}}, c_{\text{in}}}, \qquad
b^{\text{fused}}_{c_{\text{out}}} = \mathrm{scale}[c_{\text{out}}] \cdot b^{\text{conv}}_{c_{\text{out}}} + \mathrm{shift}[c_{\text{out}}].
\]
The folded conv reads identical numbers at inference but saves
the explicit BN normalization at every event. Implemented in
\texttt{hybrid\_common.py::\_fuse\_conv\_bn}.

\subsection{GPU kernel changes for the 2-layer head}

The original kernel had a single linear head:
\texttt{preds = head\_W $\cdot$ L7\_out + head\_b} with
\texttt{HEAD\_OUT = 7}. The trained head is structurally different,
not a fusable approximation.

Added to the kernel:

\begin{enumerate}[itemsep=0pt]
  \item Ten new pointer fields on \texttt{HybridS2S3Config}:
        \texttt{head\_\{cls,reg\}\_conv\_\{W,b\}},
        \texttt{head\_\{cls,reg,obj\}\_pred\_\{W,b\}}.
  \item A \texttt{\_\_device\_\_} helper \texttt{do\_two\_layer\_head}
        that runs the per-event 96$\to$96, SiLU, 96$\to$\{2,4,1\}
        sequence using one warp per event. The 96-float intermediate
        is stored in the existing per-warp \texttt{qh} smem slot,
        which is unused after the layer pipeline finishes.
  \item A dispatch in the kernel head section:
        \texttt{cfg.head\_cls\_conv\_W != 0} selects the trained path;
        otherwise the legacy single-linear path runs (used for stub
        weights).
\end{enumerate}

The Python mirror of the config struct
(\texttt{hybrid\_common.py::CHybridS2S3Config}) gained matching
\texttt{c\_uint64} fields so that the \texttt{memmove(p\_cfg,
addressof(cfg), sizeof(cfg))} call drops the right pointers into the
right slots.

\subsection{Bug: \texttt{replace\_all} and duplicated kernel sections}

The kernel has \emph{two} structurally identical head sections, one in
\texttt{k\_ssla\_s2s3\_celled\_persistent} (single-block legacy) and
one in \texttt{k\_ssla\_s2s3\_celled\_persistent\_multi} (the
production multi-block kernel that the live runner actually launches).
The two sections differ slightly only in their leading comment line.

When the head dispatch was originally added via a search-and-replace
\texttt{replace\_all=true}, only the first occurrence was matched and
updated. The second --- the production one --- continued to use the
old single-linear path that triggers only on \texttt{cfg.head\_W != 0}.
When the Python side wired up the trained head, it set
\texttt{cfg.head\_W = 0} (legacy field unused) and
\texttt{cfg.head\_cls\_conv\_W = <trained ptr>}. The kernel still ran
the old branch, saw \texttt{head\_W == 0}, and skipped the head
entirely. The pinned predictions buffer was therefore never written.
On the host side, \texttt{sigmoid(0) = 0.5}, hence the symptom: every
displayed object score was 0.5.

The diagnosis used a kernel-side \texttt{printf} of all 10 head
pointers, which confirmed that the pointers \emph{were} reaching the
GPU correctly (struct offsets matched between Python ctypes and the
CUDA struct). That left "code path not taken" as the only remaining
hypothesis, leading to the second head section.

\textbf{Generalizable lesson:} when an apparent search-and-replace
covers multiple call sites of the same logic, \emph{open the second
site and confirm} that the text matched. Identical-by-intent does not
imply identical-by-character. A trailing comment, a missing newline,
or different surrounding whitespace silently halves the replacement.

\subsection{Resolution-mismatch effect on the trained head}

The DSEC checkpoint was trained at $640 \times 430$, stride 8, giving
an s3 grid of approximately $80 \times 54$. Initial deploy testing ran
at $160 \times 120$ (25\% linear ROI on the DVXplorer's 640$\times$480
sensor), giving an s3 grid of $20 \times 15$. The architecture itself
is fully spatially translation-equivariant, but the head's
\texttt{log\_w}, \texttt{log\_h} regression was trained on absolute
pixel scales that no longer match: real objects at 25\% ROI appear
4$\times$ smaller than the trained scale.

Empirically:

\begin{tabularx}{\linewidth}{lrrr}
\toprule
\textbf{Run-time input} & \textbf{max obj\_sigmoid} & \textbf{p99 obj\_sigmoid} & \textbf{obj\_logit range} \\
\midrule
$160 \times 120$ (25\% ROI)             & 0.015 & 0.011 & $[-17, -4]$ \\
$640 \times 480$ (native, no ROI)       & 0.95  & 0.50  & $[-18, +3]$ \\
\bottomrule
\end{tabularx}

The model is essentially unable to make a confident detection at
$160 \times 120$ --- the head was never trained for that scale.

\textbf{Rule:} for visual verification of a trained checkpoint, run at
the trained input resolution. Re-cropping to a smaller window is a
training-distribution change as severe as swapping the camera.

% ====================================================================
\section{Failed Experiments (Do Not Repeat)}

The most valuable artifact of a careful optimization campaign is the
list of attempts that didn't work. Each entry below was measured
end-to-end and reverted.

\subsection{TF32 Tensor Core for L4 \texttt{qvg} projection}

Hypothesis: replace the per-event L4 \texttt{qvg} matvec
(C1 = 24 $\to$ 3 $\cdot$ C2 = 144) with TF32
\texttt{mma.sync.m16n8k4}, exploiting Ampere TC.

Implementation: a prototype that pre-packs the \texttt{qvg} matrix into
the m=16 format and dispatches the mma.sync instruction across a warp.

Result: \textbf{P1 unchanged} (\texttt{max|}$\Delta$\texttt{|} = 4.20,
TF32 precision is not a bottleneck for this network). \textbf{$-5\%$
offline throughput}, \textbf{$-3\%$ live}.

Root cause: the m=16 packing required dropping the 121 KB
\texttt{qvgIn} smem cache to fit the TC operand buffer. The smem cache
was worth more than the TC instruction.

\textbf{Do not retry} without first eliminating the smem cache
dependency (e.g.\ by re-pinning the cache in L2 instead of smem ---
itself unverified).

\subsection{Batch-local hidden-cell coalescing}

Hypothesis: within a batch, group the (event, delta) tasks by target
hidden cell, then load/store \texttt{h\_cell} once per group instead of
once per task. The motivating observation was that nearby events
often hit the same cell.

Implementation: in the layer step, sort the task queue by cell id and
fold repeats.

Result: \textbf{P1 unchanged}. \textbf{$-6\%$ throughput}.

Root cause: synthetic event distributions have reuse factor 1.097
(barely above 1). The bookkeeping overhead (sort + coalesce) was
700$\times$ the saving. The script
\texttt{orin/analyze\_coalescing.py} is retained for re-running the
analysis on real DVS camera streams where the reuse factor might
differ, but the gating heuristic must be very high before this
becomes worth the overhead.

\subsection{Raising GPU occupancy to 2 blocks/SM}

Hypothesis: doubling the number of CUDA blocks per SM (from 1 to 2)
would hide memory latency and roughly double throughput.

Measurement: 16-block layout at 96 registers/thread $\to$ 248 kev/s vs
8-block baseline 227 kev/s. Gain: \textbf{$+1.1\times$, not the
projected $+1.9\times$}.

Root cause: the kernel is compute/memory-bound at saturation, not
latency-bound. Both blocks fight for the same L2 bandwidth and the
same scheduler slots. Pipeline-split variants (separate S2 and S3
kernels, or 4-layer split kernels) hit the same wall.

\textbf{Do not retry} as the headline path to higher throughput. Path
past the current 230 kev/s ceiling likely requires:

\begin{itemize}[itemsep=0pt]
  \item a topological reorganization that reduces per-event work
        (e.g.\ trading the BATCH$=$8 cell-owner partition for a
        BATCH$=$4 with 2 blocks/SM at lower per-event smem),
  \item channel reduction in L4--L7 via distillation,
  \item or moving to a newer Jetson with more SMs.
\end{itemize}

\subsection{Single-linear head as a fused approximation of the trained head}

Hypothesis: replace the (96 $\to$ 96 BN $\to$ SiLU $\to$ 96 $\to$ 7)
trained head with a single linear (96 $\to$ 7) using the BN-fused conv
matrix composed with the prediction matrix, treating SiLU as identity.

Result: with the linearization, obj confidences during real motion
stayed in the noise floor regardless of input scale --- the SiLU
non-linearity is structural to the head's behaviour and cannot be
dropped.

Conclusion: trained-weight inference \emph{requires} a real 2-layer
head on the GPU. Not approximable.

% ====================================================================
\section{Methodology: How To Optimize Without Wasting Time}
\label{sec:methodology}

The optimization rounds documented above all followed the same
discipline. The discipline is the most portable lesson from this
campaign.

\subsection{P1 oracle is gate 0}

Every kernel change must pass the P1 oracle
(\texttt{bench\_s2\_s3\_head\_celled.py}) before any performance number
is meaningful. The oracle compares \texttt{drain\_n} output against a
single-thread CPU reference; tolerance is \texttt{max|}$\Delta$\texttt{|}
$\leq 5.0$ on \texttt{s3\_feat} and \texttt{drift = 0} on drop counts.

If P1 fails, stop benchmarking. A wrong answer at full speed is just
noise.

\subsection{Trust experiments, not projections}

Every optimization estimate in this codebase has been off by
2--10$\times$, almost always overestimating. The actual gains versus
projections:

\begin{tabularx}{\linewidth}{lXr}
\toprule
\textbf{Optimization} & & \textbf{Projection $\to$ actual} \\
\midrule
L4 \texttt{qvgIn} smem cache    & & $+20$--$30\%$ $\to$ $+4\%$ \\
L7 L2 persistence               & & $+5$--$10\%$ $\to$ $+1.7\%$ offline, $0\%$ live \\
Pipeline split / 2 blocks/SM    & & $+50$--$100\%$ $\to$ $+9\%$ then $+1.1\times$ at higher occupancy \\
Balanced topology               & & $+20$--$30\%$ $\to$ $+4\%$ \\
VG-only state-only patch        & & $+6$--$13\%$ $\to$ $+3\%$ \\
TF32 TC at L4                   & & uncertain (expert promised gain) $\to$ $-5\%$ \\
Batch-local hidden coalescing   & & uncertain (expert promised gain) $\to$ $-6\%$ \\
\bottomrule
\end{tabularx}

\textbf{Rule:} for any proposed optimization, write the smallest patch
that can deliver the gain, measure it, and only scale up if the
measurement is within 2$\times$ of the projection. Otherwise, debug or
abandon.

\subsection{One change per commit}

Bundling makes regressions hard to bisect. Each commit message records:

\begin{itemize}[itemsep=0pt]
  \item what changed (one sentence),
  \item why it should help (one paragraph of analysis),
  \item measured P1 result (drift, max$|\Delta|$),
  \item measured perf result (throughput, latency),
  \item revert decision if applicable.
\end{itemize}

\subsection{Save dead-end results}

A negative result is more valuable than a missing measurement. The
"failed experiments" section above only exists because the dead ends
were written down at the time. The auto-memory entries
\texttt{project\_tf32\_tc\_l4\_prototype}, \texttt{project\_coalescing\_net\_negative},
and \texttt{project\_occupancy\_gives\_little} encode them as machine-readable
notes that survive across conversations.

% ====================================================================
\section{Common Pitfalls and Footguns}

A grab-bag of issues that bit us during this campaign and are worth
documenting at the head of any future port.

\subsection{NVRTC PTX cache: 80--115 s rebuild after kernel edit}

The first compile of a kernel variant takes 80--115 seconds. Subsequent
runs hit the cache in $\sim$0 s. The cache lives at
\texttt{\~{}/.cache/openeva-orin-ptx/} (or
\texttt{\~{}/.cache/openeva/orin\_ptx/}) and is keyed by the hash of the
kernel source string. Editing the source --- even adding a comment ---
changes the hash and forces a rebuild.

Set \texttt{OPENEVA\_ORIN\_NO\_PTX\_CACHE=1} to skip the cache entirely
(useful when verifying cache correctness or after suspected struct
layout changes that didn't trigger an obvious rebuild).

\subsection{\texttt{\_\_launch\_bounds\_\_(288, 1)} is mandatory for the celled kernel}

NVRTC's default \texttt{MAX\_THREADS\_PER\_BLOCK} is 256. The celled
kernel uses 9 warps $\times$ 32 = 288 threads. Without the explicit
hint, \texttt{cuLaunchKernel} returns
\texttt{LAUNCH\_OUT\_OF\_RESOURCES} (error 701).

\subsection{\texttt{/tmp/ssla\_s\_*/} is volatile across reboots}

Orin clears \texttt{/tmp} on reboot. If \texttt{hybrid\_runner} fails
with "cannot open \texttt{/tmp/ssla\_s\_*/meta.json}", regenerate via
\texttt{tools/make\_ssla\_stub.py} (stub) or
\texttt{tools/convert\_dsec\_ckpt.py} (trained).

\subsection{\texttt{kept} column includes pad duplicates in P1 bench}

When the P1 oracle pads block ring sizes to a multiple of BATCH = 8,
the duplicates pass through the kernel and inflate the \texttt{kept}
count. Dedupe via global event index when comparing to the CPU
reference.

\subsection{\texttt{cv2.dnn.NMSBoxes} asserts \texttt{score\_threshold $\geq$ 0}}

Already covered in \S 6.2. The dynamic threshold for top-K fallback can
go slightly negative; clamp.

\subsection{DVXplorer Micro contrast threshold valid range is 0--17}

Already covered in \S 6.1. The getter does not reveal the silent clamp.

% ====================================================================
\section{Quick Reference}

\subsection{Build}

\begin{lstlisting}[language=Bash]
cmake -S . -B build -DCMAKE_BUILD_TYPE=Release
cmake --build build -j
\end{lstlisting}

Produces three shared libraries: \texttt{libstage01\_to\_gpu.so} (the
production hybrid pipeline), \texttt{libstage01\_capi.so} (CPU-only
baseline through \(L_3\)), \texttt{libstage0\_capi.so} (CPU-only
through \(L_1\)).

\subsection{Generate stub weights}

\begin{lstlisting}[language=Bash]
python3 tools/make_ssla_stub.py /tmp/ssla_s_64x80/ --h 64 --w 80
\end{lstlisting}

\subsection{Convert a trained checkpoint}

\begin{lstlisting}[language=Bash]
python3 tools/convert_dsec_ckpt.py \
  /path/to/best_epoch=N-step=M.ckpt \
  /tmp/ssla_s_dsec_480x640/ --h 480 --w 640
\end{lstlisting}

\subsection{P1 correctness gate}

\begin{lstlisting}[language=Bash]
cd orin
python3 bench_s2_s3_head_celled.py --n 200000
# expect: drift = 0, max|delta| <= 5
\end{lstlisting}

\subsection{Live run, synthetic source}

\begin{lstlisting}[language=Bash]
cd orin
python3 hybrid_runner_multi.py \
    --weights /tmp/ssla_s_64x80/ --h-full 64 --w-full 80 \
    --gpu-blocks 8 --shards 6 \
    --cpp-synth --synthetic-mev 2.0 --duration-s 8
\end{lstlisting}

\subsection{Live run, DVXplorer Micro, trained checkpoint}

\begin{lstlisting}[language=Bash]
cd orin
python3 hybrid_runner_multi.py \
    --weights /tmp/ssla_s_dsec_480x640/ --h-full 480 --w-full 640 \
    --gpu-blocks 8 --shards 6 \
    --roi-pct 100 --contrast 17 --cam-subsample EVERY_PIXEL \
    --demo-display --display-fps 15 --display-scale 1 \
    --obj-thresh 0.3 --nms-iou 0.65 \
    --duration-s 60
\end{lstlisting}

\subsection{Repository layout (production-relevant only)}

\begin{lstlisting}
ssla-rt/
  include/                       public C/C++ headers
  src/lib_stage01_to_gpu.cpp     production hybrid pipeline (CPU s0+s1 -> ring -> GPU)
  src/lib_stage01_capi.cpp       CPU-only stages 0+1
  src/ssla_kernels.cpp           per-stage SSLA-S forward kernels (NEON-optimized)
  tools/make_ssla_stub.py        random-weight stub generator
  tools/convert_dsec_ckpt.py     DSEC ckpt -> deploy schema conversion
  orin/orin/hybrid_common.py     ctypes struct mirrors, helpers
  orin/orin/multi_block.py       per-block resource allocation + topology
  orin/orin/weights_ssla.py      reshape ssla weights into matvec-friendly layout
  orin/kernels/ssla_s2_s3_head_celled.cuh    PRODUCTION GPU kernel
  orin/kernels/proto_layer_pair.cuh           warp-cooperative primitives
  orin/hybrid_runner_multi.py    N-block production live runner
  orin/bench_s2_s3_head_celled.py P1 oracle harness
  orin/perf_celled_multi.py      offline throughput
  docs/PIPELINE_DESIGN.md        canonical pipeline + constraints
  docs/PROFILE.md                current measurements
  docs/DEV_GUIDE.md              methodology playbook
\end{lstlisting}

% ====================================================================
\section{Closing Notes}

The system in its current state delivers $\sim$230 kev/s GPU drain and
$\sim$330\,\(\mu\)s p50 end-to-end latency at saturation on Jetson Orin
NX with stub weights, and runs the trained DSEC checkpoint end-to-end
through the production kernel via a real 2-layer YOLOX head on the GPU.

The remaining headroom is not on the CPU --- which already produces
1.85--2.0 Mev/s, well above the 230 kev/s GPU ceiling --- but on the
GPU side. The paths considered for breaking past 230 kev/s
(higher occupancy, pipeline splitting, TF32 TC) all gave $\leq 1.1\times$ in
measurement. A more invasive change (channel reduction via knowledge
distillation, or a different topology that trades BATCH for blocks per
SM) is the most plausible next direction. None of those are shipped.

For visual verification of a trained checkpoint, always run at the
trained input resolution; the head's \texttt{log\_w}, \texttt{log\_h}
regression does not transfer across resolution.

Everything in this handbook should be reproducible from the repo at
the May 2026 head-of-line commit. The authoritative numbers live in
\texttt{docs/PROFILE.md}, refreshed every measurement-bearing
optimization commit.

\end{document}
